\chapter{Reproducibility and Evaluation}
\label{ch:reproducibility}

\section{Build and Proof Commands}

The dissertation can be built from the \code{dissertation/} directory with:

\begin{lstlisting}[language=bash]
make
\end{lstlisting}

The EasyCrypt declaration inventory used in Appendix~\ref{app:ec-inventory} is
regenerated with:

\begin{lstlisting}[language=bash]
make inventory
\end{lstlisting}

This writes \code{etc/easycrypt-declaration-counts.csv}.  The target is not part
of the default PDF build because it inspects the separated proof manifest rather
than LaTeX sources.

The EasyCrypt proof gate is checked from \code{haetae-security/} with:

\begin{lstlisting}[language=bash]
EC_ATTEMPTS=8 ./provable-security/verify-provable-security.sh
\end{lstlisting}

A successful proof run ends with \code{RESULT: PASS}.

\section{Why Attempts Are Configurable}

Long EasyCrypt batches may occasionally fail because the Why3 server cannot be
started or contacted.  The verification script retries files, but accepts a file
only if EasyCrypt eventually returns success for that file.  This distinguishes
infrastructure lifecycle failures from real failed proof obligations.

\section{Evaluation Criteria}

A reviewer should check manifest completeness, manifest minimality, successful
compilation, theorem traceability, and assumption transparency.  The separated
folder is designed to support exactly that review.
